\section{问题定义、目标与系统边界}

\subsection{核心目标}
本项目的生产目标是构造一个参数化、可验证的水下无线电能传输（UWPT）电热代理系统。当前电磁模型在冻结工作频率 $\omega$ 下采用 $e^{+i\omega t}$ 峰值复相量；current-controlled 模式的 operating input 为规定的复端口电流 $c$。用户输入为几何 $g$、$c$、初始热状态与查询时间 $t$，输出包括热状态、温度、端口阻抗和功率等工程量。

系统不以黑箱方式直接学习时间轨迹。已知的 Maxwell 线性性、current-quadratic Joule structure、端口 reciprocity/passivity、矩阵级功率守恒、热质量/导热算子和连续时间动力学必须保留为硬结构或显式物理层。

若实际工作点由电压源、谐振补偿网络或负载决定，则必须在本模型输出的 $Z(a,g)$ 外接确定性的 circuit equations 求解端口电流；规定电流模型本身不等价于完整 source--load closed-loop WPT circuit。

\subsection{冻结的生产主架构}
生产主链冻结为
\begin{equation}
\boxed{
 g
 \xrightarrow{\text{validated Maxwell truth}}
 \{Z_{\rm field}(g),D_{\rm vol}(g),H_1(g),\ldots,H_r(g)\}
 \xrightarrow{\text{neural parametric surrogate}}
 \{\widehat Z_{\rm field},\widehat D_{\rm vol},\widehat H_j\}
 \xrightarrow{\text{exact current/circuit physics}}
 q(a,g,c)
 \xrightarrow{\text{true thermal ROM}}
 a(t).
}
\end{equation}
其中 $Z_{\rm field}$ 为不含独立 wire/internal impedance model 的场阻抗，$D_{\rm vol}=X^HH_\sigma X$ 为物理海水体耗散矩阵，$H_j$ 为第 $j$ 个 thermal mode 对应的 Hermitian Joule quadratic tensor。神经网络仅逼近几何到这些低维参数化物理对象的映射；端口电流不进入网络，而通过解析 quadratic forms 或外部 circuit equations 进入功率和热源。

在 reciprocal constitutive/boundary assumptions 成立时，生产解码必须保持
\begin{equation}
Z_{\rm field}^T=Z_{\rm field},
\qquad
D_{\rm vol}=D_{\rm vol}^H\succeq0,
\qquad
H_j^H=H_j,
\end{equation}
以及联合能量可行性
\begin{equation}
\operatorname{Herm}(Z_{\rm field})-D_{\rm vol}\succeq0.
\end{equation}
进一步，由 $H_j$ 与 $D_{\rm vol}$ 来自同一个非负空间 Joule density，必须满足
\begin{equation}
\boxed{
\phi_j^{\min}D_{\rm vol}
\preceq H_j
\preceq\phi_j^{\max}D_{\rm vol}.
}
\end{equation}
这些有限矩阵约束是生产可行性的必要结构，但不是 exact spatial realizability 的充分证明；完整正确性仍由 Physics Gate 和 held-out truth audit 约束。

\subsection{当前材料模型下的进一步简化}
当前正式材料模型中，线圈铜不作为 Maxwell volume conductor 的完整 solid-conductor region 处理，海水电导率未设置温度系数，介电常数和磁导率也未设置温度依赖。因此当前版本满足
\begin{equation}
\boxed{A_{\rm em}=A_{\rm em}(g),}
\end{equation}
而不是 $A_{\rm em}(a,g)$。因此固定几何的 Maxwell 场响应只需求解一次；热状态只通过明确的 conductor-temperature functional 与独立 wire impedance/resistance model 进入 wire heating 和总端口关系。

若未来引入温度相关海水/介质电磁参数，则主架构扩展为
\begin{equation}
(s_{\rm em},g)\mapsto\{Z_{\rm field},D_{\rm vol},H_j\},
\end{equation}
其中 $s_{\rm em}$ 只包含真正影响 Maxwell constitutive law 的状态变量。

\subsection{理论适用范围与生产域}
当前体系是单频、线性 time-harmonic Maxwell 与慢 thermal dynamics 的 frequency--transient 分层。cycle-averaged Joule heating 要求载波/场建立时间远短于 thermal/operating-envelope 时间尺度；quasi-static phasor circuit layer 还要求 circuit settling time 足够短。若这些尺度分离失败，应升级 EM/circuit dynamics，而不是继续沿用当前方程。

必须冻结 validated geometry domain $\mathcal G_{\rm prod}$、operating-current/circuit domain、thermal/constitutive state domain 与 initial-condition family。任何域外输入都不继承域内 validated accuracy 声明。

若未来 frequency 本身成为输入，则逐频率 passivity 不足以保证 broadband causal passivity，必须额外满足 positive-real/causality 结构；该问题不属于当前单频冻结体系。

\subsection{理论正确性的前置 Gate}
生产训练开始前，底层离散物理必须通过以下前置条件：
\begin{itemize}
  \item Maxwell formulation、reciprocity assumptions、开放域/域扩展与 mesh convergence 已验证；
  \item peak-phasor、reaction voltage、terminal/source normalization、阻抗符号和独立 Poynting/power balance 一致；
  \item $H_\sigma/W_j/D_{\rm vol}/H_j$ 使用统一的 FE/quadrature 能量定义，modal Loewner bounds 成立；
  \item PML/artificial absorption、physical seawater loss 与 wire/internal loss 严格分区，无双重计损；
  \item port current path、terminal source/sink 或 return path 满足 continuity 与明确端口定义，source self term 使用固定物理 regularization 而非 mesh regularization；
  \item 几何样本满足包封、无穿透和边界裕量要求，source/material deposition 守恒且对小几何变化足够连续；
  \item thermal basis 在构造最终 $H_j$ labels 前冻结，并覆盖完整 current-source span、initial-condition family 和多个时间尺度；rank criterion 使用 energy/dual-norm error，并通过 full transient audit；
  \item thermal coercivity/nullspace policy、pointwise temperature audit、wire-temperature functional、steady-state stability 与 circuit operating mode 均有明确定义。
\end{itemize}
这些条件没有通过时，神经网络误差不能被解释为最终物理误差。

\subsection{明确非目标}
本项目不把以下方案作为生产主路线：
\begin{enumerate}
  \item 直接学习 $(g,c,a_0,t)\mapsto T(t)$ 的黑箱 trajectory network；
  \item 用网络直接输出最终温度/Joule heat/工作点而绕过已知的 current/circuit/thermal physics；
  \item 用可行投影或 neural loss 掩盖 truth-model、port/source、boundary、wire 或 thermal-model 错误；
  \item 在底层 PDE、端口、边界、功率恒等式未通过 Physics Gate 前，以训练 loss 下降替代物理验证；
  \item 把有限测试集经验误差称为全参数域严格 certificate。
\end{enumerate}
